%-------------------------------------------------------------------------------
%  SECTION TITLE
%-------------------------------------------------------------------------------
\cvsection{\'Education}

%-------------------------------------------------------------------------------
%  CONTENT
%-------------------------------------------------------------------------------
\begin{cventries}

  %---------------------------------------------------------
  \cventry
  {LIRMM -- I3S -- CNRS AIST-JRL -- Supervisée par \emph{Dr. Andrew Comport} et \emph{Dr. Abderrahmane
  Kheddar}.} % Organization
  {Doctorat en Robotique} % Job title
  {Soutenu en Nov. 2018} % Location / Date alternative pour clarté
  {Oct. 2014 -- Nov. 2018} % Date(s)
  {
    \begin{cvitems} % Description(s) of tasks/responsibilities
    \item \entrypositionstyle{Titre} : « SLAM visuel pour la localisation et la commande en boucle fermée de
      robots humanoïdes ».
    \item \entrypositionstyle{Cadre de Recherche} - Thèse adossée à des projets d'envergure internationale
      :\emph{DARPA Robotics Challenge (DRC)}, \emph{Robohow} et \emph{H2020 COMANOID}
    \item \entrypositionstyle{Publications}: 3 publications en premier auteur -- finaliste au prix du meilleur
      article à SII 2016.
    \item \entrypositionstyle{Contributions}: Intégration d'algorithmes de \keyword{SLAM} pour la
      localisation et contrôle de
      robot humanoïdes ; méthodes de recalage entre CAO et nuage de points (\keyword{ICP}). Intégration de
      ces travaux à un algorithme de marche par commande prédictive (MPC) en boucle fermée.
    \item \entrypositionstyle{Développement Logiciel} : Codes associés aux publications (\emph{ICP},
      \emph{Robcalib}, \emph{Savitzky-Golay}, ...), contributions au framework \keyword{mc\_rtc}.
    \end{cvitems}
  }

  %---------------------------------------------------------

  %---------------------------------------------------------
  \cventry
  {Polytech Nice Sophia-Antipolis -- Université de Nice} % Institution
  {Diplôme d'Ingénieur en Informatique -- Vision, Image et Multimédia} % Degree
  {Sophia Antipolis, France} % Location
  {Sept. 2011 -- Sept. 2014} % Date(s)
  {
    \begin{cvitems} % Description(s) bullet points
      % \item \entrypositionstyle{Spécialisation} : Vision, Image et Multimédia.
    \item \entrypositionstyle{ERASMUS - Trinity College Dublin (2012-2013)} : Master de technologies
      interactives digitales (jeu vidéo)
    \item \entrypositionstyle{Stage - CRANN Institute, Dublin (2013, 3 mois)}: Conception d'un logiciel
      C++/Qt de recalage de courbes pour la microscopie à effet tunnel.
    \item \entrypositionstyle{Stage - Technische Universität München, Allemagne (2014, 6 mois)} :
      Reconnaissance de lieux par réseaux de neurones convolutionnels (CNN) supervisés par \emph{Dr. J.
      Sturm} et \emph{Prof. D. Cremers}
    \item \entrypositionstyle{Projets}: Rendu photo-réaliste de cartes SLAM dans un casque VR (\emph{Dr A.
      Comport}), jeu-vidéo 3D adapté aux personnes malvoyantes
    \end{cvitems}
  }

  %---------------------------------------------------------
  \cventry
  {Lycée Kérichen} % Institution
  {Classes Préparatoires aux Grandes Écoles (CPGE) -- MPSI} % Degree
  {Brest, France} % Location
  {Sept. 2009 -- Juin 2011} % Date(s)
  {}
\end{cventries}
